% ===== MANUALLY WRITTEN MISSING PSPEC TABLES =====

\paragraph{PSPEC Proses 1.7: Integrasi API BRIVA (Pembuatan Virtual Account)}~\\
\begin{longtable}{|p{0.2\textwidth}|p{0.7\textwidth}|}
	\caption{PSPEC Proses 1.7 -- Integrasi API BRIVA (Pembuatan Virtual Account)}
	\label{table:3.pspec17}\\
	\hline
	\textbf{Nama Proses} & Integrasi API BRIVA (Pembuatan Virtual Account) \\
	\hline
	\textbf{Deskripsi Proses} & Sistem berinteraksi dengan API BRIVA untuk membuat Virtual Account (VA) secara otomatis. \\
	\hline
	\textbf{Data Input} & Permintaan pembuatan VA dari sistem berdasarkan jenis formulir. \\
	\hline
	\textbf{Data Output} & Data Virtual Account dari BRIVA (nomor VA dan nominal pembayaran). \\
	\hline
	\textbf{Body / Keterangan} & Proses ini merupakan inti dari sistem pembayaran formulir otomatis. Ketergantungan pada API BRIVA membuat proses ini penting untuk memastikan transaksi dan verifikasi pembayaran berlangsung efisien dan akurat. \\
	\hline
\end{longtable}

\paragraph{PSPEC Proses 1.8: Pembuatan Virtual Account Manual oleh Admin Validasi}~\\
\begin{longtable}{|p{0.2\textwidth}|p{0.7\textwidth}|}
	\caption{PSPEC Proses 1.8 -- Pembuatan Virtual Account Manual oleh Admin Validasi}
	\label{table:3.pspec18}\\
	\hline
	\textbf{Nama Proses} & Pembuatan Virtual Account Manual oleh Admin Validasi \\
	\hline
	\textbf{Deskripsi Proses} & Admin Validasi dapat secara manual membuat Virtual Account (VA) jika proses generate VA otomatis melalui API BRIVA mengalami kegagalan atau sistem tidak merespons. \\
	\hline
	\textbf{Data Input} & Input biaya yang perlu dibayarkan dari Admin Validasi. \\
	\hline
	\textbf{Data Output} & Nomor Virtual Account yang digenerate secara manual. \\
	\hline
	\textbf{Body / Keterangan} & Fitur ini berfungsi sebagai mekanisme cadangan untuk menjaga kelancaran proses pembayaran ketika terjadi kendala teknis pada integrasi API BRIVA. Dengan cara ini, calon mahasiswa tetap dapat melanjutkan pembayaran tanpa hambatan. \\
	\hline
\end{longtable}

\paragraph{PSPEC Proses 2.2: Aktivasi dan Hitung VA Berdasarkan Jenis Formulir}~\\
\begin{longtable}{|p{0.2\textwidth}|p{0.7\textwidth}|}
	\caption{PSPEC Proses 2.2 -- Aktivasi dan Hitung VA Berdasarkan Jenis Formulir}
	\label{table:3.pspec22}\\
	\hline
	\textbf{Nama Proses} & Aktivasi dan Hitung VA Berdasarkan Jenis Formulir \\
	\hline
	\textbf{Deskripsi Proses} & Sistem menghitung biaya formulir berdasarkan jenis formulir dan gelombang pendaftaran, kemudian mengaktifkan Virtual Account untuk pembayaran. \\
	\hline
	\textbf{Data Input} & Jenis formulir dan gelombang pendaftaran yang dipilih CAMABA. \\
	\hline
	\textbf{Data Output} & Nomor VA BRIVA aktif beserta nominal pembayaran. \\
	\hline
	\textbf{Body / Keterangan} & Sistem mengambil data harga formulir dari konfigurasi, menghitung total biaya, lalu menginisiasi pembuatan VA. Jika VA berhasil dibuat, CAMABA akan menerima informasi pembayaran. \\
	\hline
\end{longtable}

\paragraph{PSPEC Proses 2.6: Verifikasi Pembayaran Formulir (Callback BRIVA)}~\\
\begin{longtable}{|p{0.2\textwidth}|p{0.7\textwidth}|}
	\caption{PSPEC Proses 2.6 -- Verifikasi Pembayaran Formulir (Callback BRIVA)}
	\label{table:3.pspec26}\\
	\hline
	\textbf{Nama Proses} & Verifikasi Pembayaran Formulir (Callback BRIVA) \\
	\hline
	\textbf{Deskripsi Proses} & Sistem menerima notifikasi callback dari BRIVA ketika pembayaran VA formulir berhasil dilakukan oleh CAMABA. \\
	\hline
	\textbf{Data Input} & Data callback dari API BRIVA (nomor VA, jumlah pembayaran, status transaksi). \\
	\hline
	\textbf{Data Output} & Status pembayaran diperbarui menjadi lunas di database sistem. \\
	\hline
	\textbf{Body / Keterangan} & Proses verifikasi otomatis ini memastikan bahwa status pembayaran CAMABA diperbarui secara real-time setelah transaksi berhasil, sehingga CAMABA dapat langsung melanjutkan ke tahap pengisian formulir. \\
	\hline
\end{longtable}

\paragraph{PSPEC Proses 3.5: Generate Kartu Ujian}~\\
\begin{longtable}{|p{0.2\textwidth}|p{0.7\textwidth}|}
	\caption{PSPEC Proses 3.5 -- Generate Kartu Ujian}
	\label{table:3.pspec35}\\
	\hline
	\textbf{Nama Proses} & Generate Kartu Ujian \\
	\hline
	\textbf{Deskripsi Proses} & Sistem secara otomatis menghasilkan kartu ujian untuk CAMABA setelah formulir pendaftaran berhasil disubmit. \\
	\hline
	\textbf{Data Input} & Data formulir pendaftaran CAMABA yang telah disubmit. \\
	\hline
	\textbf{Data Output} & Kartu ujian lengkap dengan QR code yang dapat diunduh. \\
	\hline
	\textbf{Body / Keterangan} & Kartu ujian berisi informasi nomor ujian, jadwal, lokasi/link ujian, foto, dan QR code. Dokumen ini wajib dibawa CAMABA saat pelaksanaan ujian untuk proses validasi. \\
	\hline
\end{longtable}

\paragraph{PSPEC Proses 4.1: Validasi Nomor Ujian}~\\
\begin{longtable}{|p{0.2\textwidth}|p{0.7\textwidth}|}
	\caption{PSPEC Proses 4.1 -- Validasi Nomor Ujian}
	\label{table:3.pspec41}\\
	\hline
	\textbf{Nama Proses} & Validasi Nomor Ujian \\
	\hline
	\textbf{Deskripsi Proses} & Sistem memvalidasi nomor ujian yang diinputkan oleh CAMABA sebelum mengakses halaman ujian. \\
	\hline
	\textbf{Data Input} & Nomor ujian yang diinput CAMABA. \\
	\hline
	\textbf{Data Output} & Konfirmasi validitas nomor ujian atau pesan error jika tidak ditemukan. \\
	\hline
	\textbf{Body / Keterangan} & Jika nomor ujian valid, sistem melanjutkan ke validasi jadwal ujian. Jika tidak valid, sistem menampilkan pesan error dan meminta CAMABA untuk memeriksa kembali nomor ujiannya. \\
	\hline
\end{longtable}

\paragraph{PSPEC Proses 4.5: Notifikasi Ujian Belum Dimulai}~\\
\begin{longtable}{|p{0.2\textwidth}|p{0.7\textwidth}|}
	\caption{PSPEC Proses 4.5 -- Notifikasi Ujian Belum Dimulai}
	\label{table:3.pspec45}\\
	\hline
	\textbf{Nama Proses} & Notifikasi Ujian Belum Dimulai \\
	\hline
	\textbf{Deskripsi Proses} & Sistem menampilkan pemberitahuan kepada CAMABA bahwa jadwal ujian belum dimulai. \\
	\hline
	\textbf{Data Input} & Status waktu ujian dari pengecekan jadwal. \\
	\hline
	\textbf{Data Output} & Pesan notifikasi bahwa ujian belum dimulai beserta informasi jadwal ujian. \\
	\hline
	\textbf{Body / Keterangan} & CAMABA akan diarahkan untuk menunggu hingga waktu ujian dimulai. Sistem dapat menampilkan countdown timer hingga jadwal ujian tiba. \\
	\hline
\end{longtable}

\paragraph{PSPEC Proses 4.9: Generate Nomor Pendaftaran dan Password Tanggal Lahir}~\\
\begin{longtable}{|p{0.2\textwidth}|p{0.7\textwidth}|}
	\caption{PSPEC Proses 4.9 -- Generate Nomor Pendaftaran dan Password Tanggal Lahir}
	\label{table:3.pspec49}\\
	\hline
	\textbf{Nama Proses} & Generate Nomor Pendaftaran dan Password Tanggal Lahir \\
	\hline
	\textbf{Deskripsi Proses} & Sistem menghasilkan nomor pendaftaran unik dan menetapkan tanggal lahir sebagai password awal untuk CAMABA yang lulus. \\
	\hline
	\textbf{Data Input} & Data kelulusan CAMABA dari proses evaluasi ujian. \\
	\hline
	\textbf{Data Output} & Nomor pendaftaran unik dan password awal (tanggal lahir) untuk proses daftar ulang. \\
	\hline
	\textbf{Body / Keterangan} & Nomor pendaftaran ini akan digunakan sebagai identitas CAMABA dalam proses daftar ulang. Password awal menggunakan tanggal lahir untuk mempermudah login pertama kali. \\
	\hline
\end{longtable}

\paragraph{PSPEC Proses 4.11: Kirim Email Hasil Ujian dan Status Kelulusan}~\\
\begin{longtable}{|p{0.2\textwidth}|p{0.7\textwidth}|}
	\caption{PSPEC Proses 4.11 -- Kirim Email Hasil Ujian dan Status Kelulusan}
	\label{table:3.pspec411}\\
	\hline
	\textbf{Nama Proses} & Kirim Email Hasil Ujian dan Status Kelulusan \\
	\hline
	\textbf{Deskripsi Proses} & Sistem mengirimkan email kepada CAMABA yang berisi hasil ujian, status kelulusan, serta instruksi langkah selanjutnya. \\
	\hline
	\textbf{Data Input} & Data hasil ujian, status kelulusan, dan informasi langkah selanjutnya. \\
	\hline
	\textbf{Data Output} & Email notifikasi terkirim ke CAMABA. \\
	\hline
	\textbf{Body / Keterangan} & Email ini memberikan informasi penting mengenai hasil seleksi secara langsung dan instruksi untuk langkah selanjutnya, seperti jadwal daftar ulang bagi yang lulus. \\
	\hline
\end{longtable}

\paragraph{PSPEC Proses 5.1: Login via Email dan Password (untuk Bebas Testing)}~\\
\begin{longtable}{|p{0.2\textwidth}|p{0.7\textwidth}|}
	\caption{PSPEC Proses 5.1 -- Login via Email dan Password (untuk Bebas Testing)}
	\label{table:3.pspec51}\\
	\hline
	\textbf{Nama Proses} & Login via Email dan Password (untuk Bebas Testing) \\
	\hline
	\textbf{Deskripsi Proses} & CAMABA jalur bebas testing melakukan login menggunakan email dan password untuk memulai proses daftar ulang. \\
	\hline
	\textbf{Data Input} & Email dan password CAMABA. \\
	\hline
	\textbf{Data Output} & Akses ke halaman daftar ulang. \\
	\hline
	\textbf{Body / Keterangan} & Proses ini khusus untuk jalur bebas testing yang mungkin memiliki alur login berbeda dari jalur testing reguler. \\
	\hline
\end{longtable}

\paragraph{PSPEC Proses 5.3: Upload Berkas Ranking SMA dan Formulir Pendaftaran}~\\
\begin{longtable}{|p{0.2\textwidth}|p{0.7\textwidth}|}
	\caption{PSPEC Proses 5.3 -- Upload Berkas Ranking SMA dan Formulir Pendaftaran}
	\label{table:3.pspec53}\\
	\hline
	\textbf{Nama Proses} & Upload Berkas Ranking SMA dan Formulir Pendaftaran \\
	\hline
	\textbf{Deskripsi Proses} & CAMABA jalur bebas testing bersyarat (ranking) mengunggah bukti ranking SMA dan formulir pendaftaran yang diperlukan. \\
	\hline
	\textbf{Data Input} & Berkas ranking SMA dan formulir pendaftaran. \\
	\hline
	\textbf{Data Output} & Berkas terunggah untuk validasi. \\
	\hline
	\textbf{Body / Keterangan} & Berkas-berkas ini akan divalidasi secara manual oleh Admin Validasi untuk memastikan kelengkapan dan keabsahannya. \\
	\hline
\end{longtable}

\paragraph{PSPEC Proses 5.5: Kirim Notifikasi Berkas Tidak Valid ke Camaba}~\\
\begin{longtable}{|p{0.2\textwidth}|p{0.7\textwidth}|}
	\caption{PSPEC Proses 5.5 -- Kirim Notifikasi Berkas Tidak Valid ke Camaba}
	\label{table:3.pspec55}\\
	\hline
	\textbf{Nama Proses} & Kirim Notifikasi Berkas Tidak Valid ke Camaba \\
	\hline
	\textbf{Deskripsi Proses} & Sistem mengirimkan notifikasi kepada CAMABA apabila berkas yang diunggah dianggap tidak valid oleh Admin Validasi. \\
	\hline
	\textbf{Data Input} & Status tidak valid dari proses validasi berkas Admin Validasi. \\
	\hline
	\textbf{Data Output} & Notifikasi terkirim ke CAMABA. \\
	\hline
	\textbf{Body / Keterangan} & Notifikasi ini memberitahu CAMABA bahwa ada masalah dengan berkas mereka dan mungkin memerlukan perbaikan atau unggah ulang berkas yang benar. \\
	\hline
\end{longtable}

\paragraph{PSPEC Proses 6.1: Validasi Data Camaba (Daftar Ulang)}~\\
\begin{longtable}{|p{0.2\textwidth}|p{0.7\textwidth}|}
	\caption{PSPEC Proses 6.1 -- Validasi Data Camaba (Daftar Ulang)}
	\label{table:3.pspec61}\\
	\hline
	\textbf{Nama Proses} & Validasi Data Camaba (Daftar Ulang) \\
	\hline
	\textbf{Deskripsi Proses} & Admin Validasi melakukan verifikasi kelengkapan dan kebenaran data CAMABA yang telah menyelesaikan daftar ulang. \\
	\hline
	\textbf{Data Input} & Data daftar ulang CAMABA. \\
	\hline
	\textbf{Data Output} & Status validasi (diterima/ditolak/perlu perbaikan). \\
	\hline
	\textbf{Body / Keterangan} & Admin Validasi dapat melihat data yang dikelompokkan per fakultas dan program studi, kemudian melakukan validasi manual berdasarkan kelengkapan berkas yang disubmit. \\
	\hline
\end{longtable}

\paragraph{PSPEC Proses 6.3: Kirim Pesan / Notifikasi oleh Admin Validasi}~\\
\begin{longtable}{|p{0.2\textwidth}|p{0.7\textwidth}|}
	\caption{PSPEC Proses 6.3 -- Kirim Pesan / Notifikasi oleh Admin Validasi}
	\label{table:3.pspec63}\\
	\hline
	\textbf{Nama Proses} & Kirim Pesan / Notifikasi oleh Admin Validasi \\
	\hline
	\textbf{Deskripsi Proses} & Admin Validasi dapat mengirim pesan atau notifikasi langsung ke email atau notifikasi di web CAMABA. \\
	\hline
	\textbf{Data Input} & Pesan atau notifikasi dari Admin Validasi. \\
	\hline
	\textbf{Data Output} & Notifikasi dan email log tercatat. \\
	\hline
	\textbf{Body / Keterangan} & Fitur ini memungkinkan Admin Validasi untuk berkomunikasi dengan CAMABA, misalnya untuk meminta perbaikan data yang ditolak atau memberikan informasi tambahan. \\
	\hline
\end{longtable}

\paragraph{PSPEC Proses 7.1: Kelola Periode Pendaftaran}~\\
\begin{longtable}{|p{0.2\textwidth}|p{0.7\textwidth}|}
	\caption{PSPEC Proses 7.1 -- Kelola Periode Pendaftaran}
	\label{table:3.pspec71}\\
	\hline
	\textbf{Nama Proses} & Kelola Periode Pendaftaran \\
	\hline
	\textbf{Deskripsi Proses} & Admin Pusat menentukan dan mengatur jadwal gelombang pendaftaran, batas waktu, dan syarat-syarat pendaftaran. \\
	\hline
	\textbf{Data Input} & Jadwal gelombang, syarat pendaftaran, dan batas waktu. \\
	\hline
	\textbf{Data Output} & Konfigurasi periode pendaftaran tersimpan di sistem. \\
	\hline
	\textbf{Body / Keterangan} & Admin Pusat memiliki kendali penuh atas periode pendaftaran. Pengaturan ini berdampak langsung pada ketersediaan layanan pendaftaran bagi CAMABA. \\
	\hline
\end{longtable}

\paragraph{PSPEC Proses 7.6: Pengumuman dan Notifikasi untuk Camaba}~\\
\begin{longtable}{|p{0.2\textwidth}|p{0.7\textwidth}|}
	\caption{PSPEC Proses 7.6 -- Pengumuman dan Notifikasi untuk Camaba}
	\label{table:3.pspec76}\\
	\hline
	\textbf{Nama Proses} & Pengumuman dan Notifikasi untuk Camaba \\
	\hline
	\textbf{Deskripsi Proses} & Admin Pusat dapat membuat dan mengirimkan pengumuman atau notifikasi kepada seluruh CAMABA terkait informasi penting PMB. \\
	\hline
	\textbf{Data Input} & Isi pengumuman atau notifikasi dari Admin Pusat. \\
	\hline
	\textbf{Data Output} & Pengumuman ditampilkan di dashboard CAMABA dan/atau notifikasi email terkirim. \\
	\hline
	\textbf{Body / Keterangan} & Fitur ini memastikan CAMABA mendapatkan informasi terkini terkait jadwal, persyaratan, atau perubahan kebijakan PMB secara tepat waktu. \\
	\hline
\end{longtable}

\paragraph{PSPEC Proses 4.10: Atur Jadwal Publikasi Kelulusan dan No Pendaftaran}~\\
\begin{longtable}{|p{0.2\textwidth}|p{0.7\textwidth}|}
\caption{PSPEC Proses 4.10 -- Atur Jadwal Publikasi Kelulusan dan No Pendaftaran}
\label{table:3.pspec410}\\
\hline
\textbf{Nama Proses} & Atur Jadwal Publikasi Kelulusan dan No Pendaftaran \\
\hline
\textbf{Deskripsi Proses} & Admin Validasi mengatur kapan hasil kelulusan dan nomor pendaftaran dapat dipublikasikan dan diakses oleh CAMABA. \\
\hline
\textbf{Data Input} & Input dari Admin Validasi. \\
\hline
\textbf{Data Output} & Jadwal publikasi kelulusan. \\
\hline
\textbf{Body / Keterangan} & Proses ini memberikan kontrol kepada Admin Validasi untuk merilis hasil kelulusan sesuai jadwal yang telah ditetapkan. \\
\hline
\end{longtable}

\paragraph{PSPEC Proses 4.10: Atur Jadwal Publikasi Kelulusan dan No Pendaftaran}~\\
\begin{longtable}{|p{0.2\textwidth}|p{0.7\textwidth}|}
\caption{PSPEC Proses 4.10 -- Atur Jadwal Publikasi Kelulusan dan No Pendaftaran}
\label{table:3.pspec410}\\
\hline
\textbf{Nama Proses} & Atur Jadwal Publikasi Kelulusan dan No Pendaftaran \\
\hline
\textbf{Deskripsi Proses} & Admin Validasi mengatur kapan hasil kelulusan dan nomor pendaftaran dapat dipublikasikan dan diakses oleh CAMABA. \\
\hline
\textbf{Data Input} & Input dari Admin Validasi. \\
\hline
\textbf{Data Output} & Jadwal publikasi kelulusan. \\
\hline
\textbf{Body / Keterangan} & Proses ini memberikan kontrol kepada Admin Validasi untuk merilis hasil kelulusan sesuai jadwal yang telah ditetapkan. \\
\hline
\end{longtable}
